\section{Threats to Validity}
\label{sec:threats}

\subsection{Construct validity}
\label{sec:threats-construct}

The four mutation operators represent a deliberately narrow subset of
reproducibility-relevant choices in ML research software. They do not cover all
sources of experimental variability, configuration drift, data-processing
choices, hardware effects, nondeterminism, or environment differences that may
affect reproducibility.

Accordingly, the mutations should be interpreted as controlled
reproducibility-relevant perturbations rather than as a complete model of real
ML faults. In particular, the study does not assume that every applied mutation
corresponds to a naturally occurring defect or that all reproducibility
failures can be represented by the four operators.

This limitation is partly intentional. The supported mutations were restricted
to mechanically identifiable syntax with explicit transformations, making
candidate selection and mutation application auditable. The resulting gain in
measurement precision comes at the cost of narrower construct coverage.

A second construct threat concerns the validation workflow used as the
detection oracle. The selected workflow represents one frozen repository
validation path, not the complete set of safeguards that may exist in a
repository. A SURVIVED mutation therefore means only that the selected workflow
did not detect the applied confirmed non-equivalent change. It does not imply
that no other test, experiment, manual inspection, or downstream scientific
analysis would detect it.

Similarly, a KILLED mutation establishes sensitivity to one applied mutation
but does not establish that the repository is generally protected against
reproducibility failures.

We mitigate these interpretation risks by separating workflow type from oracle
type, explicitly distinguishing completion-only workflows from workflows with
substantive acceptance conditions, and treating mutation outcomes as properties
of repository--mutation--workflow combinations rather than repository-level
reproducibility labels.


\subsection{Internal validity}
\label{sec:threats-internal}

A central internal-validity risk in mutation studies is outcome-dependent
selection. Repository choice, mutation-candidate choice, validation-workflow
choice, or oracle classification could otherwise be influenced by knowledge of
which cases are likely to produce KILLED or SURVIVED outcomes.

For the primary B02 corpus, we addressed this risk through static,
outcome-blind corpus construction. Repository eligibility, fixed revisions,
candidate mutations, validation workflows, and prospectively recorded oracle
categories were frozen before primary mutation execution. Candidate
repositories were not executed during corpus construction, and the corpus was
not extended or replaced after observing empirical outcomes.

The RQ2 classification was likewise separated from outcome analysis. Workflow
categories were inherited from the frozen corpus, and B02 oracle categories had
been recorded prospectively under protocol version 2. The RQ2 classification
frame was independently frozen before outcome joining. Consequently, the
observed RQ2 pattern was not used to redefine oracle categories.

Execution and restoration introduce another internal-validity concern. Changes
made solely to make historical software executable could alter the behavior
being measured. D027 therefore preserved the frozen repository revision,
mutation candidate, validation workflow, and oracle classification. Source
compatibility patches were prohibited, validation workflows could not be
weakened or skipped, restoration was baseline-first, and the restoration recipe
was fixed before evaluating the corresponding mutant. Baseline and mutant were
evaluated under symmetric restored conditions.

Nevertheless, restoration cannot eliminate all uncertainty associated with
executing historical software in a contemporary environment. A reconstructed
environment may differ from the original authors' environment in ways that are
not fully observable from repository metadata. We therefore retain the
canonical primary result separately and treat D027 as an additional bounded
evidence layer rather than as a replacement for primary execution.

Semantic-equivalence handling also affects internal validity. A surviving
syntactic mutation that does not change the relevant semantics should not
contribute evidence of weak validation. We therefore exclude confirmed
equivalent mutations from the meaningful detection denominator and distinguish
confirmed non-equivalent, confirmed equivalent, unverified, and not-run
semantic states.


\subsection{External validity}
\label{sec:threats-external}

The study does not constitute a random sample of all ML research repositories.
The corpus was constructed using operator-specific static eligibility criteria
and a fixed stopping rule. Consequently, the results should not be interpreted
as population-wide estimates of mutation detection or research-software
reproducibility.

The B02 operator frames are also unequal. The frozen stopping rule yielded ten
cases each for \texttt{random-seed} and \texttt{dependency-pin}, but the
eligible static frames for \texttt{data-split} and \texttt{cv-fold-count} were
exhausted at six and three cases. The meaningful evaluated denominators are
smaller still.

The repository population may also differ systematically from other forms of ML
software. The study focuses on research repositories for which both a supported
mutation candidate and a usable repository validation workflow could be
identified under the static protocol. Production ML systems, notebooks without
repository-level workflows, closed-source research code, and projects using
unsupported frameworks or configuration mechanisms are therefore outside the
observed frame.

These constraints limit generalization, but they also define the scope of the
empirical claim precisely: the study measures the behavior of the selected
frozen repository--operator cases under their selected validation workflows.


\subsection{Conclusion validity}
\label{sec:threats-conclusion}

The number of detected mutations is small. Across the 23 confirmed
non-equivalent evaluated mutations, only two were KILLED. RQ2 category sizes
are also sparse and uneven; for example, the prospective B02 oracle analysis
contains only two meaningful \texttt{assertion} cases.

We therefore treat RQ2 as descriptive and exploratory. We do not infer causal
effects of workflow type or oracle type, and we do not interpret differences in
observed percentages as evidence of statistical superiority.

This is especially important for the B02 oracle comparison. Both detected
mutations occurred in the
\texttt{dependency-pin} $\times$ \texttt{completion-only} cell. The observed
difference between \texttt{assertion} and \texttt{completion-only} workflows
can therefore reflect operator composition, sparse denominators, differential
evaluability, or other case-level factors rather than oracle strength.

For the same reason, operator-specific proportions such as the observed
detection among confirmed non-equivalent \texttt{dependency-pin} cases are
reported descriptively and should not be interpreted as estimates of a general
operator effect.

The primary quantitative claim is therefore intentionally narrow: in this
sample, the selected validation workflows detected 2 of 23 confirmed
non-equivalent reproducibility-relevant mutations.


\subsection{Temporal and executability validity}
\label{sec:threats-temporal}

The repositories in the corpus were executed after their frozen revisions had
been created, using contemporary software and infrastructure. Historical
dependency ecosystems, package indexes, compilers, operating-system libraries,
hardware availability, and external resources may have changed in the
intervening period.

This temporal distance was empirically consequential. Primary execution
evaluated only 13 of 39 frozen cases, and bounded restoration increased the
combined evaluated set to 24. The remaining non-evaluable cases therefore
reduce the observable empirical denominator and may introduce differential
attrition across workflow, operator, or oracle categories.

We do not interpret present-day execution failure as direct evidence that the
original research artifact was irreproducible at the time of publication.
Instead, these failures characterize the difficulty of retrospectively
executing the frozen software under the study conditions.

Conversely, successful contemporary restoration does not reconstruct the
original environment perfectly. D027 was designed to recover bounded
executability while preserving the frozen empirical target, not to claim
historical environment identity.

The study therefore reports primary and restored evidence separately before
combining successful evaluations at the analysis layer. This separation makes
the effect of temporal software-ecosystem drift visible rather than silently
removing non-evaluable cases from the study.
